\section{Introducción}

La evaluación ambiental de la bioconversión con \textit{Hermetia illucens} no debe limitarse a balances estáticos o factores de emisión promedio. En sistemas biológicos altamente dinámicos, las emisiones cambian en función del tiempo, de la humedad del sustrato, de la temperatura y del estado metabólico del sistema. Por ello, la incorporación de métricas dinámicas permite una valoración más realista del desempeño ambiental del proceso.

Este artículo analiza la utilidad de un enfoque de evaluación ambiental dinámico para interpretar las emisiones de CO$_2$ y CH$_4$ durante la bioconversión. La propuesta se centra en traducir las salidas del sistema híbrido de monitoreo en indicadores útiles para la gestión, la comparación tecnológica y la toma de decisiones operativas. De esta manera, el análisis supera la lógica de una estimación estática y avanza hacia una lectura temporal del comportamiento ambiental del proceso.

\section{Enfoque de evaluación}

El enfoque parte de la idea de que una medición promedio puede ocultar episodios críticos de emisión y, por tanto, subestimar riesgos ambientales relevantes. En contraste, una aproximación dinámica permite identificar ventanas temporales de mayor intensidad metabólica, acumulación de carbono emitido y eventos asociados a condiciones de anaerobiosis. Esta información resulta especialmente valiosa en sistemas BSF, donde la interacción entre larvas, sustrato y microbiota genera trayectorias no lineales de emisión.

A partir de estas series temporales es posible construir métricas ambientales con mayor utilidad aplicada, tales como emisiones acumuladas, tasas instantáneas, duración de eventos críticos y comparación entre escenarios operativos. Estas métricas permiten contrastar la interpretación obtenida mediante factores estáticos frente a la obtenida con monitoreo continuo. Así, el artículo propone una base conceptual para un ACV dinámico o una evaluación ambiental temporalmente explícita del sistema de bioconversión.

\section{Implicaciones para la gestión}

Desde la perspectiva de gestión ambiental, el principal valor del enfoque dinámico es que transforma las emisiones en señales de control del proceso. En lugar de usarse solo como resultado final, las emisiones pueden servir para detectar ineficiencias, ajustar condiciones operativas y priorizar intervenciones sobre aireación, humedad o carga orgánica. Esto mejora la capacidad de respuesta frente a episodios que comprometen tanto el desempeño ambiental como la estabilidad biológica del sistema.

Además, el uso de métricas dinámicas fortalece la comparabilidad entre tratamientos y mejora la trazabilidad requerida en contextos de reporte, verificación y evaluación de sostenibilidad. En este sentido, el artículo aporta una aproximación metodológica para vincular monitoreo inteligente, evaluación ambiental y toma de decisiones en procesos de valorización de residuos con BSF.
